% !TeX root = ieice_draft.tex
%% Manuscript for IEICE Technical Report.
%% Compile with: latexmk -pdf ieice_draft.tex
\documentclass{ieice-tecrep}
\usepackage[T1]{fontenc}
\usepackage[utf8]{inputenc}
\usepackage{graphicx}
\usepackage[fleqn]{amsmath}
\usepackage{newtxtext}
\usepackage[varg]{newtxmath}
\usepackage{booktabs}
\usepackage{etoolbox}
\usepackage{cite}
\usepackage{needspace}
\usepackage[hidelinks]{hyperref}
\newtheorem{proposition}{Proposition}
\newtheorem{lemma}{Lemma}
\newtheorem{theorem}{Theorem}
% Keep the class's default list spacing for Pandoc-origin lists.
\providecommand{\tightlist}{}
\setlength{\emergencystretch}{1em}
\tecrepyear{2026}
\title{Extending the L-star Algorithm to Learn Syntactic Monoids}
\subtitle{}
% Use commas between authors; override only this manuscript's output macro.
\makeatletter
% Keep theorem headings and statements flush left, with wrapping titles.
\newcommand{\draft@theoremlist}{%
 \Needspace{5\baselineskip}%
 \list{}{\leftmargin\z@\rightmargin\z@\labelwidth\z@\labelsep\z@
  \itemindent\z@\listparindent\z@
  \topsep.5\Cvs\parsep\z@\partopsep\z@\itemsep\z@}%
 \item[]\phantomsection}
\renewcommand{\@begintheorem}[2]{%
 \draft@theoremlist\def\@currentlabelname{#1}%
 \textbf{\@thmlbracket #1\ #2\@thmrbracket}\hspace{1em}\ignorespaces}
\renewcommand{\@opargbegintheorem}[3]{%
 \draft@theoremlist\def\@currentlabelname{#3}%
 \textbf{\@thmlbracket #1\ #2\@thmrbracket\ (#3)}\hspace{1em}\ignorespaces}
\renewcommand{\@endtheorem}{\endlist}
\patchcmd{\@outputEauthorlist}{\ifnum\@totalnumber=\tw@\ \else , \fi}{}{}{\PackageError{ieice-draft}{Author separator patch failed}{Check the class version.}}
\patchcmd{\@outputEauthorlist}{and\ \nobreak}{,\ \nobreak}{}{\PackageError{ieice-draft}{Author separator patch failed}{Check the class version.}}
\makeatother
% Both authors share the affiliation label, following readme.tex.
\DeclareRobustCommand{\NagoyaAddress}{%
 Furo-cho, Chikusa-ku, Nagoya%
}
\authorlist{%
 \authorentry[zhanght@sqlab.jp]{Haotian Zhang}{Nagoya}
 \authorentry[y-inoue@sqlab.jp]{Yusuke Inoue}{Nagoya}
}
\affiliate[Nagoya]{Nagoya University\\\NagoyaAddress}
\begin{document}
\begin{abstract}
This essay introduces an algorithm for learning
the syntactic monoid of an unknown regular language based on the original L-star algorithm. 
We will define the extended observation table and the
corresponding hypothesis acceptor, explain essential properties including closedness, right
consistency, and monoid consistency, and analyze the complexity of the
learning process. By monoid consistency, we will introduce two strategies to learn syntactic monoids.
Future work will implement both strategies in a simulated active
learning environment, examine which classes of regular languages
favor each strategy, and compare their query and computational costs
with learning a minimal DFA using \(L^*\) and then constructing its
transition monoid.
\end{abstract}
\begin{keyword}
syntactic monoid, regular language, active automata learning, observation table, monoid consistency
\end{keyword}
\maketitle

\section{Introduction}\label{introduction}

Angluin's \(L^*\) algorithm actively learns the minimal deterministic
finite automaton (DFA) of an unknown regular language through
membership and equivalence queries~\cite{angluin87}. In this paper,
\(U\subseteq\Sigma^*\) denotes the unknown regular language, and
\(\Sigma\) denotes a fixed finite alphabet throughout. The learner
does not require access to an internal representation of \(U\):
membership queries determine whether a word belongs to \(U\), while
equivalence queries test a hypothesis and return a counterexample when
it is incorrect. Active automata learning has also been used beyond
the original setting; for example, Khmelnitsky et al.\ learn a DFA
surrogate of a recurrent neural network for property-directed
verification~\cite{khmelnitsky2020property}.
Vazquez-Chanlatte et al.\ combine active automata learning with large
language models to learn DFA task specifications from natural language,
labeled examples, and expert demonstrations~\cite{vazquezchanlatte2025lstarlm}.
Chen et al.\ introduce CEDAR, which applies active automata learning
to construct DFAs for embodied skills and natural-language constraints,
then composes them to enforce the learned constraints during
execution~\cite{chen2026cedar}.

A minimal DFA captures the behavior of a language through its states
and transitions. The syntactic monoid additionally exposes the algebra
of word concatenation, making algebraic properties of the language
accessible. For example, a regular language is star-free exactly when
its syntactic monoid is aperiodic~\cite{SCHUTZENBERGER1965190}.
The syntactic monoid is isomorphic to the transition monoid of the
minimal DFA. A conventional route to learning it is therefore to run
the original \(L^*\) algorithm to obtain the minimal DFA and then
compute the transformations of its state set induced by all words.
These transformations form the DFA's transition monoid, which is
isomorphic to the syntactic monoid. Our aim is to construct a
representation of the syntactic monoid within the active learning
process itself.

These two representations arise from different equivalence relations
on words. The states of the minimal DFA correspond to classes of the
Myhill--Nerode right congruence, which distinguishes words by their
possible suffixes. In contrast, the syntactic monoid is the quotient
of \(\Sigma^*\) by syntactic congruence, which distinguishes words by
both left and right contexts. Its right Cayley graph describes
multiplication by the input letters. Taking the identity element as
the initial state and marking the classes contained in \(U\) as
accepting turns this graph into a DFA recognizing \(U\).

We extend \(L^*\) to learn this algebraic representation. Our
extended observation table records both left and right contexts, and
a monoid consistency condition checks whether words assigned the same
row behave identically in the contexts needed for syntactic
congruence. We give two strategies for maintaining this condition,
\texttt{post\_check} and \texttt{Promote\_to\_P}. Both use membership
and equivalence queries to the teacher; the monoid consistency check
uses information obtained through membership queries rather than a
third kind of teacher query.

The remainder of this paper is organized as follows.
Section~\ref{preliminaries} introduces the necessary terminology.
Section~\ref{algorithm} describes the extended observation table and
the two learning strategies. Section~\ref{main-results} establishes
their correctness, Section~\ref{complexity-analysis} analyzes their
complexity, and Section~\ref{conclusion} concludes the paper.

\section{Preliminaries}\label{preliminaries}

\subsection{Words and Languages}\label{minimal-dfa-preliminaries}

Let \(\Sigma\) be a finite nonempty alphabet. A \emph{word} over \(\Sigma\)
is a finite sequence of letters from \(\Sigma\). The set of all such
words is denoted by \(\Sigma^*\), and \(\epsilon\) denotes the empty
word. The length of a word \(w\) is written as \(|w|\), with
\(|\epsilon|=0\). The concatenation of words \(u\) and \(v\) is
written as \(uv\) or \(u\cdot v\). A \emph{language} over \(\Sigma\) is a
subset of \(\Sigma^*\). A language is \emph{regular} if it is accepted by a
deterministic finite automaton, as defined below.

A word \(u \in \Sigma^*\) is a \emph{prefix} of \(w \in \Sigma^*\) if \(w=uv\) for some
\(v\in\Sigma^*\). A set \(X\subseteq\Sigma^*\) is \emph{prefix-closed}
if every prefix of every word in \(X\) also belongs to \(X\).
Similarly, \(v \in \Sigma^*\) is a \emph{suffix} of \(w\) if \(w=uv\) for some
\(u\in\Sigma^*\), and \(X\) is \emph{suffix-closed} if it contains every
suffix of each of its words. In particular, every nonempty
prefix-closed or suffix-closed set contains \(\epsilon\).

A \emph{deterministic finite automaton (DFA)} is a tuple
\(\mathcal{A}=(Q,\Sigma,\delta,q_0,F)\), where \(Q\) is a finite
nonempty set of states, \(\delta:Q\times\Sigma\to Q\) is a total
transition function, \(q_0\in Q\) is the initial state, and
\(F\subseteq Q\) is the set of accepting states. The transition
function extends to \(\delta^*:Q\times\Sigma^*\to Q\) by
\[
\begin{aligned}
\delta^*(q,\epsilon)&=q,\\
\delta^*(q,wa)&=\delta(\delta^*(q,w),a)
\quad(w\in\Sigma^*,\ a\in\Sigma).
\end{aligned}
\]
The language accepted by \(\mathcal{A}\) is defined as 
\(L(\mathcal{A})=\{w\in\Sigma^*\mid\delta^*(q_0,w)\in F\}\).

A \emph{minimal DFA} for a regular language \(L\) is a DFA with the fewest states that accepts \(L\).
Every regular language has a unique minimal DFA \(\mathcal{A}_L\), and \(\mathcal{A}_L\) can be constructed as described follows: 
% Every state of a minimal DFA is reachable from its
% initial state, and any two distinct states \(q,q'\) can be
% distinguished by a word \(z\in\Sigma^*\): exactly one of
% \(\delta^*(q,z)\) and \(\delta^*(q',z)\) belongs to \(F\).
% Therefore, 
Let \(\equiv^{\mathrm{MN}} \subseteq \Sigma^* \times \Sigma^*\) be the equivalence relation such that
\[
w\equiv_L^{\mathrm{MN}}w^{\prime}
\Longleftrightarrow
\bigl(\forall v\in\Sigma^*,\quad wv\in L\Longleftrightarrow w'v\in L\bigr)
\] for each \(w,w'\in\Sigma^*\).
We call this relation the \emph{Myhill--Nerode right congruence} of \(L\).
This relation is a right congruence: if
\(w\equiv_L^{\mathrm{MN}}w'\), then
\(wx\equiv_L^{\mathrm{MN}}w'x\) for every \(x\in\Sigma^*\).
Write \([w]_{\mathrm{MN}}\) for the class of \(w\).
The Myhill--Nerode theorem \cite{myhill1957finite,nerode1958linear} states that \(L\) is regular if and only if
this relation has finitely many equivalence classes. These classes
form the state set of the minimal DFA: the initial state is
\([\epsilon]_{\mathrm{MN}}\), the transition on a letter \(a\) maps
\([w]_{\mathrm{MN}}\) to \([wa]_{\mathrm{MN}}\), and a class is
accepting precisely when its words belong to \(L\).

\subsection{Syntactic Monoids}\label{syntactic-monoid}

For a language \(L\subseteq\Sigma^*\), let \(\sim_L \subseteq \Sigma^* \times \Sigma^*\) be the equivalence relation such that
\[
w\sim_L w^{\prime}
\Longleftrightarrow
\bigl(\forall u,v\in\Sigma^*,\quad uwv\in L\Longleftrightarrow uw'v\in L\bigr)
\] for each \(w,w'\in\Sigma^*\).
We call this relation the \emph{syntactic congruence} of \(L\).
Thus syntactic congruence tests both left and right contexts.
It is a two-sided congruence: \(w\sim_L w'\) implies
\(uwv\sim_L uw'v\) for all \(u,v\in\Sigma^*\).
Taking the left context to be \(\epsilon\) shows that
\(w\sim_L w'\) implies \(w\equiv_L^{\mathrm{MN}}w'\).

Let \([w]_L\) denote the syntactic congruence class of \(w\).
The quotient \(M(L)=\Sigma^*/{\sim_L}\), equipped with the
multiplication \([u]_L\cdot[v]_L=[uv]_L\), is the \emph{syntactic monoid}
of \(L\). Multiplication is well-defined because \(\sim_L\) is a
two-sided congruence, and its identity is
\(1_L=[\epsilon]_L\). The canonical quotient morphism
\(\eta_L:\Sigma^*\to M(L)\) is given by \(\eta_L(w)=[w]_L\).

Let \(F_L=\{[w]_L\mid w\in L\}\). Since syntactically congruent
words have the same membership in \(L\), we have \(L=\eta_L^{-1}(F_L)\). Hence \(M(L)\) recognizes \(L\).
For a regular language \(L\), the syntactic monoid is finite. It is the
smallest monoid recognizing \(L\) and is isomorphic to the
transition monoid of the minimal DFA for \(L\)~\cite{holzer2004deterministic}.

Cayley introduced the graphical representation of groups in
1878~\cite{cayley1878graphical}. Here we use the corresponding
construction for monoids based on right multiplication.
The \emph{right Cayley graph} of \(M(L)\), with letters of \(\Sigma\)
acting as generators through \(\eta_L\), has vertex set \(M(L)\)
and a labeled edge
\(m\xrightarrow{a}m\cdot\eta_L(a)\) for each \(m\in M(L)\)
and \(a\in\Sigma\). Thus each vertex has exactly one outgoing
edge for each letter. Define the corresponding transition function
by \(\delta_L(m,a)=m\cdot\eta_L(a)\).

For regular \(L\), selecting \(1_L\) as the initial state and
\(F_L\) as the accepting states turns this graph into the DFA
\(\mathcal{M}_L=(M(L),\Sigma,\delta_L,1_L,F_L)\).
Indeed, \(\delta_L^*(1_L,w)=\eta_L(w)=[w]_L\), so a word
\(w\) is accepted exactly when \(w\in L\).
Consequently, the right Cayley graph of the syntactic monoid,
equipped with these initial and accepting states, is a DFA
recognizing \(L\). The state set of DFA \(\mathcal{M}_L\) has \(|M(L)|\) elements and is larger than that of the minimal DFA generally. 

\section{Algorithm}\label{algorithm}

In this section, we will introduce the extended \(L^*\) algorithm for learning the syntactic monoid of an unknown regular language.
First, we will introduce the concepts and properties used in the algorithm, and then we will describe the process of two different learning strategies.

\subsection{Queries and Extended Observation Table}\label{algorithm-concepts}
\textbf{Membership Queries.} In a membership query, the learner selects a word \(w \in \Sigma^*\), and the teacher answers whether or not \(w \in U\).
We denote the answer by the characteristic function
\(\chi:\Sigma^*\to\{0,1\}\), defined by
\[
\chi(w)=
\begin{cases}
1, & \text{if } w\in U,\\
0, & \text{if } w\notin U.
\end{cases}
\]
Thus a membership query accepts \(w\) when \(\chi(w)=1\) and
rejects it when \(\chi(w)=0\).

\textbf{Equivalence Queries.} When the observation table \(OT\) is closed, we can obtain a hypothesis
automaton \(\mathcal{M}\) according to the definition of the acceptor.
Then the teacher answers whether or not the unknown regular language
\(U\) is the language accepted by \(\mathcal{M}\).

If the teacher answers ``yes'' to the equivalence query, then the
language \(L(\mathcal{M})\) accepted by the hypothesis automaton
\(\mathcal{M}\) is the same as the unknown regular set \(U\).
Consequently,
\[
\forall w \in \Sigma^*,\quad
w \in L(\mathcal{M}) \Longleftrightarrow w \in U.
\]

If the teacher answers ``no'' to the equivalence query, then the teacher provides a counterexample \(w \in \Sigma^*\) such that 
\[
w \in L(\mathcal{M}) \text{ and } w \notin U.
\]
or vice versa. The learner can refine the observation table using this counterexample and repeat the learning process. 

\textbf{Extended Observation Table.} In Angluin's original \(L^*\) algorithm~\cite{angluin87}, an
observation table is a triple \((R,E,T_0)\), where \(R\) is a
nonempty finite set of candidate representative words and \(E\) is a
nonempty finite set of test words. The function
\(T_0:(R\cup R\Sigma)E\to\{0,1\}\) records membership through
\(T_0(w)=\chi(w)\) on its domain.
Rows are indexed by \(R\cup R\Sigma\), columns by \(E\), and
the entry at row \(r\) and column \(e\) is \(T_0(re)\). 

To learn syntactic congruence classes, we extend the observation table
to include left contexts as well. Our extended observation table
is a 4-tuple
\[
OT = (I,P,S,T),
\]

where  \(I,P,S \subseteq \Sigma^*\) are finite sets, and \(T\) is the function 
\[
T : P\times( I \cup I \cdot \Sigma) \times S \rightarrow \{0,1\}
\]
defined by \(T(p,i,s)=\chi(pis)\). The roles of the sets are as follows:

\begin{itemize}
\tightlist
\item
  Elements of \(I\) are candidates for right Cayley graph states of the syntactic monoid of \(U\). 
  In particular, if
  \[
  \forall p\in P,\forall s\in S,\quad
  T(p,i,s) = T(p,i',s)
  \]
  for \(i,i'\in I\cup I\Sigma\), 
  then \(i\) and \(i'\) are regarded as the same state.
\item  
  Elements of \(P\) correspond to test prefixes for these states.
\item
  Elements of \(S\) correspond to test suffixes for these states.
\item
  Elements of \(I \cdot \Sigma\) are used to construct the transition
  function.
\end{itemize}

When \(P=\{\epsilon\}\), the entries reduce
to the suffix-test table entries \(T(\epsilon,i,s)\) used in an ordinary
observation table. The additional test prefixes allow us to compare
words through the values of the form \(\chi(pis)\). The observation table is initially set to \((\{\epsilon\},\{\epsilon\},\{\epsilon\},T)\).

\begin{table}[!ht]
\centering
\caption{Structure of the extended observation table.}
\label{tab:observation-table}
\begin{tabular}{@{}ll@{}}
\toprule
Row set & Entries indexed by \(P \times S\) \\
\midrule
\(I\) & \(P \cdot I \cdot S\) \\
\(I \cdot \Sigma\) & \(P \cdot I \cdot \Sigma \cdot S\) \\
\bottomrule
\end{tabular}
\end{table}

\textbf{Rows and Status Equivalence.}
For each \(i\in I\cup I\Sigma\), define the row function
\(\mathrm{row}_{P,S}(i):P\times S\to\{0,1\}\) by
\(\mathrm{row}_{P,S}(i)(p,s)=T(p,i,s)\).
For an arbitrary word \(x\in\Sigma^*\), extend this notation by
\(\mathrm{row}_{P,S}(x)(p,s)=\chi(pxs)\). This agrees with the
stored table rows whenever \(x\in I\cup I\Sigma\), and also
defines rows for words not currently stored in the table.

We define \emph{status equivalence relation} \(\equiv_{\mathrm{status}} \subseteq \Sigma^* \times \Sigma^*\) as
\[
i \equiv_{\mathrm{status}} i'
\Longleftrightarrow
\text{row}_{P,S}(i)=\text{row}_{P,S}(i')
\]
for each \(i,i^{\prime} \in \Sigma^*\).
We use \([i]_{\mathrm{status}}\) to denote the equivalence class of
\(i\) under this status equivalence. Thus, if \(i \equiv_{\mathrm{status}} i'\), then \([i]_{\mathrm{status}} = [i']_{\mathrm{status}}\). 
Our algorithm aims that the status equivalence is equal to the syntactic congruence when the algorithm terminates. 

\textbf{Partial Right Row and Suffix Equivalence.}
We can consider only the suffix part 
\[
\text{row}_{S}(i) \in \{0,1\}^{S}
\]
defined by
\[
\text{row}_{S}(i) = \bigl(T(\epsilon,i,s)\bigr)_{s\in S}
\]
for each \(i \in I\cup I\Sigma\).
For arbitrary \(x\in\Sigma^*\), the corresponding extension is
\(\text{row}_S(x)=(\chi(xs))_{s\in S}\). 

We define \emph{suffix equivalence relation} \(\equiv_{S} \subseteq \Sigma^* \times \Sigma^*\) as 
\[
i \equiv_{S} i'
\Longleftrightarrow
\text{row}_{S}(i)=\text{row}_{S}(i')
\]
for each \(i,i' \in \Sigma^*\). In our algorithm, 
the suffix equivalence is equal to the Myhill-Nerode congruence when the algorithm terminates 
which is similar to the original \(L^*\) algorithm. 

Unless explicitly stated otherwise, \(\text{row}(x)\)
denotes the full row \(\text{row}_{P,S}(x)\). When only the suffix part
with the empty left context is used, we write \(\text{row}_S(x)\).

\textbf{Hypothesis Acceptor.}
Given an observation table \(OT = (I,P,S,T)\), the \emph{hypothesis acceptor} of \((I,P,S,T)\) is a 5-tuple \(\mathcal{M} = (Q,\Sigma,\delta,q_0,F)\), where
\[
Q = \{\text{row}(i)\mid i \in I\},
\]
\[
q_0 = \text{row}(\epsilon) \in Q,
\]
\[
F = \{\text{row}(i) \in Q \mid T(\epsilon,i,\epsilon)=1\},
\]

and \(\delta:Q \times \Sigma\rightarrow Q\) is the transition function defined as follows:
For \(q \in Q\) and \(a\in \Sigma\), if \(q = \text{row}(i)\) for some
\(i \in I\) and \( ia \equiv_{\mathrm{status}} i'\),\(i'\in I\), then \(\delta(q,a)=\text{row}(i')\). 
This function is well-defined when the observation table satisfies some properties discussed later.

If the transition function \(\delta\) is well-defined, it can be extended to \(\delta^{*}:Q \times \Sigma^{*} \rightarrow Q\) by 
\[
\delta^{*}(q,w) =
\left\{
\begin{array}{ll}
q, & \text{if } w = \epsilon, \\
\delta(\delta^*(q,w'),a), & \text{if } w = w'a,\ a \in \Sigma.
\end{array}
\right.
\]
\subsection{Table Properties and Refinement}\label{algorithm-properties-refinement}

In this part, we will introduce some properties, essential to obtain the hypothesis acceptor from our algorithm. 
In addition, we will discuss how to refine the observation table when these properties are not satisfied. 

\textbf{Separability.}
The observation table is \emph{separable} if no two distinct words in \(I\) are status-equivalent:
\[
\forall i_1,i_2 \in I,\quad
i_1 \neq i_2 \Longleftrightarrow i_1 \not\equiv_{\mathrm{status}} i_2.
\]

The process in our algorithm guarantees that the observation table is separable because 
we add a word to \(I\) only if its row differs
from every existing row in \(I\). Refining \(P\) or \(S\)
preserves separability, since adding test words cannot make distinct
rows equal. 

\textbf{Closedness.}
The observation table is \emph{closed} if for every \(i \in I\) and \(a\in \Sigma\), 
there exists \(i'\in I\) such that \(ia \equiv_{\mathrm{status}} i'\).
That is,
\[
\forall i \in I,\forall a \in \Sigma,\exists i' \in I,\quad
ia \equiv_{\mathrm{status}} i'.
\]

If the table is not closed, find a word
\(t \in I\cdot\Sigma\) whose row differs from all rows in \(I\)
and add \(t\) to \(I\). Fill the new table entries using membership
queries and repeat the closedness check. This preserves separability;
see Proposition~\ref{proposition-1-closedness}.

\textbf{Right Consistency.}
The observation table is \emph{right consistent} if for all \(i,i'\in I\) and \(a\in \Sigma\),
\[
i \equiv_S i'
\Longrightarrow
ia \equiv_S i'a.
\]
That is,
\[
\forall i,i' \in I,\forall a \in \Sigma,\quad
i \equiv_S i' \rightarrow ia \equiv_S i'a.
\]

If right consistency fails, find \(i,i'\in I\),
\(a\in\Sigma\), and \(s\in S\) such that \(i\equiv_S i'\) but
\(T(\epsilon,ia,s)\neq T(\epsilon,i'a,s)\). Add \(as\) to \(S\)
and fill the new table entries using membership queries. This makes
\(i\) and \(i'\) distinguishable by the suffix \(as\); then recheck
closedness and right consistency. See Proposition~\ref{proposition-2.1-right-consistency-refinement} for details.

If the observation table is separable, closed, and right consistent, then we can
define a hypothesis automaton \(\mathcal{M}\).

\textbf{Monoid Consistency.}
The observation table is \emph{monoid consistent} if for all \(v,i\in I\), \(t\in I\Sigma\), and \(s\in S\),
\[
t\equiv_{\mathrm{status}}v\Rightarrow \chi(its)=\chi(ivs).
\]

Here the left context \(i\) ranges over \(I\) and need not belong
to \(P\), so these values are not necessarily entries of the current
table. If \(i\in P\), they are precisely \(T(i,t,s)\) and
\(T(i,v,s)\), respectively.

If monoid consistency fails, find
\(t\in I\cdot\Sigma\), \(v\in I\), \(i\in I\), and \(s\in S\)
such that \(t\equiv_{\mathrm{status}}v\) but
\(\chi(its)\neq \chi(ivs)\); see Theorem~\ref{theorem-monoid-consistency-check-ii}.
Add this \(i\) to \(P\) and fill the new table entries using membership
queries. The new left context distinguishes \(t\) from \(v\);
then recheck closedness, right consistency, and monoid consistency.

\subsection{The \texttt{post\_check} Strategy}\label{strategy-post-check}

This strategy first obtains a hypothesis accepting \(U\), and then
checks whether its representatives satisfy monoid consistency.
The set \(P\) is refined only when a distinguishing left context is
needed by this check.

\textbf{Phase 1: Initialization.}
Initialize \(I=P=S=\{\epsilon\}\), and fill the observation table
using membership queries.

\textbf{Phase 2: Table Filling and Consistency Refinement.}
Check every row in \(I\cdot\Sigma\). If a word \(t\in I\cdot\Sigma\)
has a row different from all rows in \(I\), add \(t\) to \(I\) and
fill the new table entries. Repeat until the table is closed.
Also check right consistency: if \(i\equiv_S i'\) but
\(ia\not\equiv_S i'a\), find \(s\in S\) distinguishing \(ia\)
and \(i'a\), and add \(as\) to \(S\). Refill the table and repeat
the checks until it is both closed and right consistent.
New representatives are added only when their rows are distinct,
so separability is preserved.

\textbf{Phase 3: Equivalence Query and Counterexample Processing.}
Construct the hypothesis acceptor \(\mathcal{M}\) and submit it as an
equivalence query. If a counterexample \(w\) is returned, refine \(S\)
using Angluin's method~\cite{angluin87} or the Rivest--Schapire binary
search method~\cite{rivest1993inference}, and return to Phase~2. The original Angluin's method adds all prefixes of counterexample \(w\) into 
\(S\) while the Rivest-Schapire method adds only one approriate suffix of \(w\) into \(S\) which is more efficient.
If \(L(\mathcal{M})=U\), proceed to Phase~4.

\textbf{Phase 4: Monoid Consistency Check and Prefix Refinement.}
For each \(t\in I\cdot\Sigma\), find its representative \(v\in I\)
such that \(t\equiv_{\mathrm{status}}v\). Check all \(i\in I\) and
\(s\in S\) for a violation
\[
\chi(its)\neq \chi(ivs).
\]
If a violation is found, add its left context \(i\) to \(P\).
This distinguishes \(t\) from \(v\). Fill the new table entries and
return to Phase~2. If every pair passes the check, return the
hypothesis acceptor representing the learned syntactic monoid.

\subsection{The \texttt{Promote\_to\_P} Strategy}\label{strategy-promote-to-p}

This strategy uses state representatives as left-context tests from
the moment they are introduced. It maintains the invariant
\(I\subseteq P\), so left contexts required for monoid consistency
are already included in the observation table.

\textbf{Phase 1: Initialization.}
Initialize \(I=P=S=\{\epsilon\}\), and fill the observation table
using membership queries.

\textbf{Phase 2: Closedness Refinement and Promotion to \(P\).}
If a word \(t\in I\cdot\Sigma\) has a row different from all rows
in \(I\), add it to both sets:
\[
I\leftarrow I\cup\{t\},\qquad P\leftarrow P\cup\{t\}.
\]
Fill the new rows and left-context entries using membership queries,
and repeat until the table is closed. Every addition to \(I\) is
accompanied by the same addition to \(P\), preserving \(I\subseteq P\).
Separability is preserved because each new representative has a
previously unseen row.

\textbf{Phase 3: Right Consistency Refinement.}
Check the suffix rows of representatives in \(I\). If
\(i\equiv_S i'\) but \(ia\not\equiv_S i'a\), find \(s\in S\)
with \(T(\epsilon,ia,s)\neq T(\epsilon,i'a,s)\), and add \(as\)
to \(S\). Refill the table and return to Phase~2, promoting any new
representatives to \(P\) as before. Repeat until closedness and right
consistency both hold.

\textbf{Phase 4: Equivalence Query and Counterexample Processing.}
Construct \(\mathcal{M}\) and submit an equivalence query. If a
counterexample is returned, refine \(S\) using the same method as in
the \texttt{post\_check} strategy and return to Phase~2. If
\(L(\mathcal{M})=U\), return the hypothesis without a separate monoid
consistency check.

Indeed, when \(t\equiv_{\mathrm{status}}v\), equality of their full
rows already gives \(T(i,t,s)=T(i,v,s)\) for
every \(i\in I\) and \(s\in S\), because \(I\subseteq P\).
Thus the monoid consistency condition is satisfied by construction.
Compared with \texttt{post\_check}, this strategy includes more left
contexts in the table, but avoids the separate check after a successful
equivalence query.

\section{Correctness of Syntactic Monoid Learning}\label{main-results}

\subsection{Representative Function}\label{definition-representative-function}

Let \(\rho:\Sigma^*\to I\) be the path representative function defined
by the acceptor: \(\rho(\epsilon)=\epsilon\). If \(x=ya\) with \(a\in\Sigma\) and \(\rho(y)=i\in I\), then by
closedness there exists \(u\in I\) such that \(ia\equiv_{\mathrm{status}}u\). We define \(\rho(ya)=u\). Intuitively, if we input string \(x\) into the acceptor, it will get
into a state \(q = \delta^*(q_0,x)\). By the definition of the acceptor,
there is a representative \(i \in I\) which fulfills
\(q = \text{row}(i)\). We say \(i\) is the representative of \(x\) in
this acceptor.

\Needspace{8\baselineskip}
\subsection{Propositions, Lemma and Theorem}\label{propositions-lemma-and-theorem}
\begin{proposition}[Closedness]\label{proposition-1-closedness}
If \((I,P,S)\) is separable but not closed, then
using membership queries one can find \(i' \in I \cdot \Sigma\) such
that \((I \cup \{i'\}, P,S)\) remains separable.
\end{proposition}
\noindent\textbf{Proof.} By the definition of closedness, there exists
\(i' \in I \cdot \Sigma\) such that \(\forall i \in I,\quad i \not\equiv_{\mathrm{status}} i'\). We add this \(i'\) into \(I\) so that separability is preserved.
\hfill $\square$

\begin{proposition}[Right Consistency Refinement]\label{proposition-2.1-right-consistency-refinement}
If \((I,P,S)\) is separable and closed but not
right consistent, then we can refine \(S\) so that these two locally
equivalent words become distinguishable.
\end{proposition}
\noindent\textbf{Proof.} According to the definition of right consistency, there
exist \(i,i' \in I\) and \(a\in \Sigma\) such that \(i\equiv_S i' \quad\text{and}\quad ia \not\equiv_S i'a\). In other words, for such \(i\), \(i'\), and \(a\),
\[
\exists s \in S,\quad
T(\epsilon,ia,s) \neq T(\epsilon,i'a,s).
\]

When we add \(as\) into \(S\), the observation table checks
\(\text{row}(i)\) and \(\text{row}(i')\) using \(T(\epsilon,i,as) \neq T(\epsilon,i',as)\). Thus, \(i \not\equiv_S i'\). Therefore, refining \(S\) by adding \(as\) makes these two words
distinguishable. \hfill $\square$

\begin{proposition}[Equivalence Queries and Suffix-Equivalent Representatives]\label{proposition-2.2-status-equivalence-implies-suffix-equivalence}
Assume that the observation table
\(OT=(I,P,S,T)\) is closed, separable, and right consistent.
Let \(\mathcal{M}\) be the hypothesis acceptor constructed from \(OT\),
and let \(\rho:\Sigma^*\to I\) be its path representative function.
Assume that \(\mathcal{M}\) passes the equivalence query, that is, \(L(\mathcal{M})=U\). Then for every \(w\in\Sigma^*\), \(w\equiv_S\rho(w)\).
\end{proposition}
\noindent\textbf{Proof.} Since \(I\) is prefix-closed, following the prefixes of
any \(i\in I\) in the acceptor gives
\(\delta^*(q_0,i)=\text{row}(i)\). By the definition of the path
representative function,
\[
\delta^*(q_0,w)=\text{row}(\rho(w))
=\delta^*(q_0,\rho(w)).
\]

Therefore, for every \(s\in S\),
\[
\begin{aligned}
\delta^*(q_0,ws)
&=\delta^*(\delta^*(q_0,w),s)\\
&=\delta^*(\delta^*(q_0,\rho(w)),s)\\
&=\delta^*(q_0,\rho(w)s).
\end{aligned}
\]

Thus \(ws\) and \(\rho(w)s\) are either both accepted or both rejected
by \(\mathcal{M}\). Since \(L(\mathcal{M})=U\), we obtain
\[
\forall s\in S,\quad ws\in U\Longleftrightarrow\rho(w)s\in U.
\]

By the definition of suffix equivalence, this means
\(w\equiv_S\rho(w)\).
\par\nopagebreak\noindent\hfill $\square$

\begin{proposition}[Right Consistency and Myhill-Nerode Congruence]\label{proposition-2.3-right-consistency-and-myhill-nerode-congruence}
Assume that the observation table
\(OT=(I,P,S,T)\) is closed, separable, and right consistent.
Let \(\mathcal{M}\) be the hypothesis acceptor constructed from \(OT\),
and let \(\rho:\Sigma^*\to I\) be its path representative function.
Assume that \(\mathcal{M}\) passes the equivalence query, that is, \(L(\mathcal{M})=U\). Then for all \(u,v\in\Sigma^*\),
\[
u\equiv_Sv
\Rightarrow
\bigl(\forall y\in\Sigma^*,\quad uy\in U\Longleftrightarrow vy\in U\bigr).
\]
\end{proposition}
\noindent\textbf{Proof.} First, let \(i,i'\in I\) satisfy
\(i\equiv_S i'\). We prove by induction on \(|y|\) that
\(\rho(iy)\equiv_S\rho(i'y)\) for every \(y\in\Sigma^*\).

\noindent\textbf{Base case.} If \(y=\epsilon\), then
\(\rho(i)=i\) and \(\rho(i')=i'\), because \(I\) is prefix-closed
and the table is separable. Thus the claim follows from
\(i\equiv_S i'\).

\noindent\textbf{Induction step.} Assume that
\(\rho(iy)\equiv_S\rho(i'y)\), and let \(a\in\Sigma\).
Since both representatives belong to \(I\), right consistency gives
\(\rho(iy)a\equiv_S\rho(i'y)a\).
By the definition of \(\rho\),
\[
\begin{aligned}
\rho(iy)a&\equiv_{\mathrm{status}}\rho(iya),\\
\rho(i'y)a&\equiv_{\mathrm{status}}\rho(i'ya).
\end{aligned}
\]
Since \(\epsilon\in P\), status equivalence implies suffix
equivalence. Therefore,
\[
\rho(iya)\equiv_S\rho(iy)a
\equiv_S\rho(i'y)a\equiv_S\rho(i'ya).
\]
This completes the induction.

Since \(\epsilon\in S\), the representatives \(\rho(iy)\) and
\(\rho(i'y)\) have the same accepting status. By the definition of
the path representative function, these representatives correspond
to the states reached after reading \(iy\) and \(i'y\), respectively.
Hence \(iy\in L(\mathcal M)\) if and only if
\(i'y\in L(\mathcal M)\). As \(L(\mathcal M)=U\), we obtain
\(iy\in U\) if and only if \(i'y\in U\).

For any strings \(u,v\in\Sigma^*\), assume that \(u\equiv_S v\).
Proposition~\ref{proposition-2.2-status-equivalence-implies-suffix-equivalence}
gives \(u\equiv_S\rho(u)\) and \(v\equiv_S\rho(v)\); hence
\(\rho(u)\equiv_S\rho(v)\). Since these representatives belong to
\(I\), the preceding induction shows that
\(\rho(u)y\in U\) if and only if \(\rho(v)y\in U\) for every
\(y\in\Sigma^*\). Furthermore, by the definition of the path
representative function,
\(\delta^*(q_0,u)=\delta^*(q_0,\rho(u))\) and
\(\delta^*(q_0,v)=\delta^*(q_0,\rho(v))\). Thus \(u\) and
\(\rho(u)\) have the same accepting behavior after every suffix,
as do \(v\) and \(\rho(v)\). Since \(L(\mathcal M)=U\), for every
\(y\in\Sigma^*\) we obtain
\[
uy\in U \Longleftrightarrow \rho(u)y\in U \Longleftrightarrow \rho(v)y\in U \Longleftrightarrow vy\in U.
\]
\hfill $\square$

\begin{proposition}[Right Compatibility of Status Equivalence]\label{proposition-2.4-stable-right-transition-of-status-equivalence-i-dont-have-any-idea-about-the-name-of-this-title}
Assume that the observation table
\(OT=(I,P,S,T)\) is closed, separable, and right consistent.
Let \(\mathcal{M}\) be the hypothesis acceptor constructed from \(OT\),
and let \(\rho:\Sigma^*\to I\) be its path representative function.
Assume that \(\mathcal{M}\) passes the equivalence query, that is, \(L(\mathcal{M})=U\). Then for all \(u,v\in\Sigma^*\) and \(a\in\Sigma\),
\[
u\equiv_{\mathrm{status}}v
\Rightarrow
ua\equiv_{\mathrm{status}}va.
\]
\end{proposition}
\noindent\textbf{Proof.} Assume \(u\equiv_{\mathrm{status}}v\). Then for every
\(p\in P\) and \(s\in S\), \(\chi(pus)=\chi(pvs)\). Thus, for every \(p\in P\), \(pu\equiv_S pv\). By Proposition~\ref{proposition-2.3-right-consistency-and-myhill-nerode-congruence}, for every \(z\in\Sigma^*\), \(puz\in U\Longleftrightarrow pvz\in U\). Take \(z=as\). Then for every \(p\in P\) and \(s\in S\), \(puas\in U\Longleftrightarrow pvas\in U\). Therefore, \(ua\equiv_{\mathrm{status}}va\).
\hfill $\square$


\begin{proposition}[Status equivalence and Representative]\label{proposition-2.5-status-equivalence-and Representative}
Assume that the observation table
\(OT=(I,P,S,T)\) is closed, separable, and right consistent.
Let \(\mathcal{M}\) be the hypothesis acceptor constructed from \(OT\),
and let \(\rho:\Sigma^*\to I\) be its path representative function.
Assume that \(\mathcal{M}\) passes the equivalence query, that is, \(L(\mathcal{M})=U\). Then for every \(x\in\Sigma^*\), \(x\equiv_{\mathrm{status}}\rho(x)\). Equivalently, \(\operatorname{row}(x)=\operatorname{row}(\rho(x))\).
\end{proposition}
\noindent\textbf{Proof.} We prove the statement by induction on \(|x|\).

\textbf{Base case.} If \(x=\epsilon\), then by definition \(\rho(\epsilon)=\epsilon\). Therefore, \(\epsilon\equiv_{\mathrm{status}}\rho(\epsilon)\).

\textbf{Induction step.} Let \(x=ya\), where \(a\in\Sigma\). Assume that \(y\equiv_{\mathrm{status}}\rho(y)\). Let \(\rho(y)=i\). By Proposition~\ref{proposition-2.4-stable-right-transition-of-status-equivalence-i-dont-have-any-idea-about-the-name-of-this-title}, since \(\equiv_{\mathrm{status}}\) is right
compatible, we have \(ya\equiv_{\mathrm{status}}ia\). By the definition of \(\rho\), we choose \(\rho(ya)\in I\) such that \(ia\equiv_{\mathrm{status}}\rho(ya)\). By transitivity of \(\equiv_{\mathrm{status}}\), we obtain \(ya\equiv_{\mathrm{status}}\rho(ya)\). Thus, \(x\equiv_{\mathrm{status}}\rho(x)\). This completes the induction. Therefore, for every \(x\in\Sigma^*\), \(x\equiv_{\mathrm{status}}\rho(x)\). This proves the proposition.
\hfill $\square$
% Proposition~\ref{proposition-2.5-status-equivalence-and Representative} shows that the status equivalence of any string \(x\) is the same as that of its representative \(\rho(x)\) in \(I\). This means that the representative function \(\rho\) preserves the status equivalence classes, and thus the rows of \(x\) and \(\rho(x)\) are identical.

\begin{proposition}[Status Equivalence and transition function]\label{proposition-2.6-status-equivalence-and-transition-function}
Assume that the observation table
\(OT=(I,P,S,T)\) is closed, separable, and right consistent.
Let \(\mathcal{M}\) be the hypothesis acceptor constructed from \(OT\),
and let \(\rho:\Sigma^*\to I\) be its path representative function.
Assume that \(\mathcal{M}\) passes the equivalence query, that is, \(L(\mathcal{M})=U\). Then for all \(x,y\in\Sigma^*\),
\[
x\equiv_{\mathrm{status}}y
\Longleftrightarrow
\delta^*(q_0,x)=\delta^*(q_0,y).
\]
\end{proposition}
\noindent\textbf{Proof.} Let \(\rho:\Sigma^*\to I\) be the path representative function. That is,
for every \(x\in\Sigma^*\), \(\delta^*(q_0,x)=q_{\rho(x)}\). By Proposition~\ref{proposition-2.5-status-equivalence-and Representative}, for every \(x\in\Sigma^*\), \(x\equiv_{\mathrm{status}}\rho(x)\). First, assume that \(x\equiv_{\mathrm{status}}y\). By Proposition~\ref{proposition-2.5-status-equivalence-and Representative}, we have \(x\equiv_{\mathrm{status}}\rho(x)\) and \(y\equiv_{\mathrm{status}}\rho(y)\). Therefore, by transitivity of \(\equiv_{\mathrm{status}}\), \(\rho(x)\equiv_{\mathrm{status}}\rho(y)\). Since \(\rho(x),\rho(y)\in I\) and the observation table is separable,
we get \(\rho(x)=\rho(y)\). Hence,
\[
\delta^*(q_0,x)=q_{\rho(x)}=q_{\rho(y)}=\delta^*(q_0,y).
\]

Conversely, assume that \(\delta^*(q_0,x)=\delta^*(q_0,y)\). By the definition of \(\rho\), this means \(q_{\rho(x)}=q_{\rho(y)}\). Since states are represented by elements of \(I\), we have \(\rho(x)=\rho(y)\). By Proposition~\ref{proposition-2.5-status-equivalence-and Representative}, \(x\equiv_{\mathrm{status}}\rho(x)\) and \(y\equiv_{\mathrm{status}}\rho(y)\). Because \(\rho(x)=\rho(y)\), we obtain \(x\equiv_{\mathrm{status}}y\). Therefore,
\[
x\equiv_{\mathrm{status}}y
\Longleftrightarrow
\delta^*(q_0,x)=\delta^*(q_0,y).
\]

This proves the proposition. \hfill $\square$

\begin{proposition}[Well-definedness of the Hypothesis Automaton]\label{proposition-3-well-definedness-of-the-hypothesis-automaton}
If the observation table \(OT=(I,P,S,T)\) is closed and separable, then the transition function \(\delta:Q\times\Sigma\to Q\) of the hypothesis automaton \(\mathcal{M}\) is well-defined.
Consequently, for all \(q\in Q\) and \(w\in\Sigma^*\), the extended
transition \(\delta^*(q,w)\) exists and is unique.
\end{proposition}
\noindent\textbf{Proof.} First, we prove existence. By the definition of the
acceptor, for \(i\in I\) and \(a\in\Sigma\), if there exists \(i'\in I\)
such that \(ia\equiv_{\mathrm{status}}i'\), then \(\delta(\text{row}(i),a)=\text{row}(i')\). Because of closedness, we have
\[
\forall i\in I,\forall a\in\Sigma,\exists i'\in I,\quad
ia\equiv_{\mathrm{status}}i'.
\]

Since \(i'\in I\), we have \(\text{row}(i')\in Q\). Therefore, \(\delta(\text{row}(i),a)\) exists.

Now we prove uniqueness. Suppose that there are two representatives
\(i_1',i_2'\in I\) such that \(ia\equiv_{\mathrm{status}}i_1'\) and \(ia\equiv_{\mathrm{status}}i_2'\). Then \(i_1'\equiv_{\mathrm{status}}i_2'\). Because the observation table is separable and \(i_1',i_2'\in I\), we
get \(i_1'=i_2'\). Therefore, the representative of \(ia\) in \(I\) is unique, and hence \(\delta(\text{row}(i),a)\) exists and is unique.

Since every state \(q\in Q\) has the form \(q=\text{row}(i)\) for some \(i\in I\), the above argument shows that for every \(q\in Q\)
and \(a\in\Sigma\), \(\delta(q,a)\) exists and is unique.

Finally, we prove that the extended transition function \(\delta^*\) is
also well-defined. We use induction on \(|w|\).

If \(|w|=0\), then \(w=\epsilon\), and by definition, \(\delta^*(q,\epsilon)=q\), which exists and is unique.

Assume that for every word \(w'\in\Sigma^*\) with \(|w'|\le k\),
\(\delta^*(q,w')\) exists and is unique. Let \(w=w'a, \qquad a\in\Sigma\). By the induction hypothesis, \(\delta^*(q,w')\) exists and is unique. Since \(\delta\) is well-defined, \(\delta(\delta^*(q,w'),a)\) also exists and is unique. Therefore, \(\delta^*(q,w)=\delta(\delta^*(q,w'),a)\) exists and is unique.

Thus, for all \(q\in Q\) and \(w\in\Sigma^*\), \(\delta^*(q,w)\) exists
and is unique. \hfill $\square$

\begin{lemma}[Monoid Consistency Check I]\label{lemma-monoid-consistency-check-i}
Assume that the observation table \(OT=(I,P,S,T)\) is closed,
separable, and right consistent. Let \(\mathcal{M}\) be the
hypothesis acceptor constructed from \(OT\), and let
\(\rho:\Sigma^*\to I\) be its path representative function.
Assume that \(\mathcal{M}\) passes the equivalence query, that is,
\(L(\mathcal{M})=U\). If \(\mathcal{M}\) is not isomorphic to the
right Cayley graph of the syntactic monoid of \(U\), equipped with
its initial and accepting states, then
\[
\exists t \in I \cdot \Sigma,\exists v \in I,\quad
t \equiv_{\mathrm{status}} v \land t \not\sim_U v.
\]
\end{lemma}
\noindent\textbf{Proof.} Assume that
\[
\forall t \in I\cdot\Sigma,\forall v\in I,\quad
(t\equiv_{\mathrm{status}}v \Rightarrow t\sim_U v)
\tag{*}
\]

We claim that
\[
\forall x\in\Sigma^*,\quad x\sim_U\rho(x)
\]
where \(\rho:\Sigma^*\to I\) is the path representative function. We prove this claim by induction on \(x\).

For the base case, if \(x=\epsilon\), then \(\rho(\epsilon)=\epsilon\),
so \(\epsilon\sim_U\rho(\epsilon)\). For the induction step, let \(x=ya\) with \(a\in\Sigma\). Let \(i=\rho(y)\). By the induction hypothesis, \(y\sim_U i\). Since syntactic congruence is a right congruence, we have \(ya\sim_U ia\). By the definition of \(\rho\), there exists \(u=\rho(ya)\in I\) such
that \(ia\equiv_{\mathrm{status}}u\). Since \(ia\in I\cdot\Sigma\) and \(u\in I\), by \((*)\) we get \(ia\sim_U u\). Therefore, \(ya \sim_U u=\rho(ya)\). Thus every word \(x\) is syntactically congruent to its representative
\(\rho(x)\in I\).

By Proposition~\ref{proposition-2.5-status-equivalence-and Representative}, \(x \equiv_{status} \rho(x)\). And because the observation table is separable, \(\rho(x)\) is the
unique representative for \(x\).

Conversely, syntactically congruent words have identical observation rows by definition. 
Hence, \(\equiv_{\mathrm{status}}\) coincides with \(\sim_U\).

Therefore, if \((*)\) held:
\[
\forall x,y \in \Sigma^*. x \equiv_{status} y \Rightarrow x \sim_U y.
\]

In other words: \(\equiv_{\mathrm{status}}\text{is equal to}\sim_U\). Hence \((*)\) is false. Therefore,
\[
\exists t\in I\cdot\Sigma,\exists v\in I,\quad
t\equiv_{\mathrm{status}}v\land t\not\sim_U v.
\]

This proves the lemma. \hfill $\square$

\begin{theorem}[Monoid Consistency Check II]\label{theorem-monoid-consistency-check-ii}
Assume that the observation table
\(OT=(I,P,S,T)\) is closed, separable, and right consistent.
Let \(\mathcal{M}\) be its hypothesis acceptor. Assume that
\(\mathcal{M}\) passes the equivalence query, that is,
\(L(\mathcal{M})=U\), but \(\mathcal{M}\) is not the right Cayley graph
of the syntactic monoid of \(U\). According to Lemma~\ref{lemma-monoid-consistency-check-i}, let
\(v \in I\) and \(t \in I\cdot \Sigma\) be such that \(t \equiv_{\mathrm{status}} v \quad\text{and}\quad t \not\sim_L v\). Then for this pair \((v,t)\),
\[
\exists i \in I,\exists s \in S,\quad
\chi(its) \neq \chi(ivs).
\]

Equivalently,
\[
(i \cdot t \cdot s \in U \land i \cdot v \cdot s \notin U)
\vee
(i \cdot t \cdot s \notin U \land i \cdot v \cdot s \in U).
\]
\end{theorem}
\noindent\textbf{Proof.} Since \(t\not\sim_L v\), there exist
\(x,y\in\Sigma^*\) that distinguish \(t\) from \(v\).
Without loss of generality, assume that \(xty\in U\) and
\(xvy\notin U\).

Let \(i=\rho(x)\in I\) be the path representative of \(x\).
By Proposition~\ref{proposition-2.2-status-equivalence-implies-suffix-equivalence}, \(x\equiv_S i\). Applying Proposition~\ref{proposition-2.3-right-consistency-and-myhill-nerode-congruence} gives
\[
\forall z\in\Sigma^*,\quad xz\in U\Longleftrightarrow iz\in U.
\]
Taking \(z=ty\) and \(z=vy\), respectively, we obtain
\(ity\in U\) and \(ivy\notin U\).

Suppose, for a contradiction, that \(it\equiv_S iv\).
By Proposition~\ref{proposition-2.3-right-consistency-and-myhill-nerode-congruence}, this would imply
\[
\forall z\in\Sigma^*,\quad itz\in U\Longleftrightarrow ivz\in U.
\]
In particular, taking \(z=y\) contradicts \(ity\in U\) and
\(ivy\notin U\). Hence \(it\not\equiv_S iv\).

By the definition of suffix equivalence, there exists \(s\in S\)
such that \(its\) and \(ivs\) have different membership-query results.
Equivalently,
\[
\chi(its)\neq \chi(ivs).
\]
Thus \(i\in I\) and \(s\in S\) provide the required witness.
\hfill $\square$

\subsection{Correctness}\label{correctness}

We now combine the preceding results to establish the correctness of
the algorithm's output. Specifically, we show that the conditions
satisfied at termination ensure that the learned quotient is
isomorphic to the syntactic monoid of the target language.
Before we introduce the proof process of the correctness of the extension algorithm, 
we need 2 prepared lemmas and show their proof processes.

\begin{lemma}[Extension for each transition step from status equivalence to syntactic congruence]\label{Extension for each transition step from status equivalence to syntactic congruence}
  Let \(OT=(I,P,S,T)\) be an observation table that
is closed, separable, right consistent, and monoid consistent. Let
\(\mathcal M\) be the acceptor constructed from \(OT\). If
\(\mathcal M\) passes the equivalence query, i.e., \(L(\mathcal M)=U\), then the equivalence relation \(\equiv_{\mathrm{status}}\) coincides
with the syntactic congruence \(\sim_L\):\[
   \forall v\in I, t\in I \cdot\Sigma. t \equiv_{status} v \Rightarrow t \sim_L v
\]
\end{lemma}
\textbf{Proof.} According to the definition of monoid consistency:
\[
\forall v \in I, t\in I \cdot \Sigma. t \equiv_{status} v \Rightarrow (\forall i \in I,s \in S. \chi(its) = \chi(ivs))
\]

Therefore \(\forall i \in I. it \equiv_S iv\). By Proposition~\ref{proposition-2.3-right-consistency-and-myhill-nerode-congruence}, \[
\forall w\in\Sigma^*,\quad
itw\in U\Longleftrightarrow ivw\in U.
\]

Consider the set \(Prev\): \[
   Prev = \{p\in \Sigma^*|\delta^*(q_0,p) = q_i\}
\]
where \(i \in I\) and \(\delta^*(q_0,i) = q_i\). In other words, \(i\)
is the representative of \(q_i\) and \(Prev\) is the set of all strings that lead to the same state \(q_i\) in the acceptor \(\mathcal M\).

By Proposition~\ref{proposition-2.2-status-equivalence-implies-suffix-equivalence}, \(\forall p \in Prev. p \equiv_S i\). 
By Proposition~\ref{proposition-2.3-right-consistency-and-myhill-nerode-congruence},
\[
   \forall z \in \Sigma^*. pz\in U \Leftrightarrow iz\in U
\]

Let \(z_1 = tw\) and \(z_2 = vw\):
\[
ptw\in U
\Longleftrightarrow
itw\in U
\Longleftrightarrow
ivw\in U
\Longleftrightarrow
pvw\in U
\]

Because \(p\) and \(w\) are any strings:
\[
   \forall v\in I, t\in I \cdot\Sigma. t \equiv_{status} v \Rightarrow t \sim_L v
\]
\hfill $\square$

\begin{lemma}[Extension for any strings from status equivalence to syntactic congruence]\label{Extension for any strings from status equivalence to syntactic congruence}
Let \(OT=(I,P,S,T)\) be an observation table that
is closed, separable, right consistent, and monoid consistent. Let
\(\mathcal M\) be the acceptor constructed from \(OT\). If
\(\mathcal M\) passes the equivalence query, i.e., \(L(\mathcal M)=U\), then:
\[
   \forall x \in \Sigma^*. x \sim_L \rho(x)
\]
where \(\sim_L\) is the syntactic congruence and \(\rho:\Sigma^*\to I\) is the path representative function. 
\end{lemma}
\textbf{Proof.} 
We prove this claim by induction on \(x\).

For the base case, if \(x=\epsilon\), then \(\rho(\epsilon)=\epsilon\),
so \(\epsilon\sim_L\rho(\epsilon)\). 
For the induction step, let \(x=ya\) with \(a\in\Sigma\). 
By the induction hypothesis, \(y \sim_L \rho(y)\) where \(\rho(y) \in I\) because of the definition of \(\rho\). 
Since syntactic congruence is a right congruence, we have \(ya\sim_L \rho(y)a\). 
By the definition of \(\rho\), \(\rho(y)a\equiv_{status}\rho(ya)\in I\). By closedness, \(\rho(ya)\in I\). 
By Lemma~\ref{Extension for each transition step from status equivalence to syntactic congruence}, \(\rho(y)a \sim_L \rho(ya)\). 
Therefore, \(ya \sim_L \rho(ya)\), in other words, \(x \sim_L \rho(x)\). This completes the induction.

Thus every word \(x\) is syntactically congruent to its representative \(\rho(x)\in I\). 
\hfill $\square$

\begin{theorem}[Correctness of extended \(L^*\) algorithm for \texttt{Post\_check} Strategy]\label{theorem-correctness-of-extended-lstar-algorithm-for-post-check-strategy}
Let \(OT=(I,P,S,T)\) be an observation table that
is closed, separable, right consistent, and monoid consistent. Let
\(\mathcal M\) be the acceptor constructed from \(OT\). If
\(\mathcal M\) passes the equivalence query, i.e., \(L(\mathcal M)=U\), then the equivalence relation \(\equiv_{\mathrm{status}}\) coincides
with the syntactic congruence \(\sim_L\): \[
\forall x,y\in\Sigma^*,\quad
x\equiv_{\mathrm{status}}y
\Longleftrightarrow
x\sim_L y.
\]
\end{theorem}
Consequently, the quotient \(\Sigma^*/{\equiv_{\mathrm{status}}}\) is isomorphic to the syntactic monoid of \(U\), and \(\mathcal M\) is
isomorphic to its right Cayley graph.

Under these assumptions, the implication
\[
   \forall x,y\in\Sigma^*,\quad
   x\sim_L y
   \Rightarrow
   x\equiv_{\mathrm{status}}y
\]
follows directly from the definition of observation rows. It therefore
suffices to prove the reverse implication:
\[
\forall x,y\in\Sigma^*,\quad
x\equiv_{\mathrm{status}}y
\Rightarrow
x\sim_L y. \tag{**}
\]
\textbf{Proof.} 
Let \(x,y\in\Sigma^*\) and assume that \(x\equiv_{\mathrm{status}}y\). 
By Proposition~\ref{proposition-2.5-status-equivalence-and Representative}, 
we have \(x\equiv_{\mathrm{status}}\rho(x)\) and \(y\equiv_{\mathrm{status}}\rho(y)\). 
Therefore, by transitivity of \(\equiv_{\mathrm{status}}\), we get \(\rho(x)\equiv_{\mathrm{status}}\rho(y)\). 
Since \(\rho(x),\rho(y)\in I\), by closedness and separability, \(\rho(x) = \rho(y)\). 
By Lemma~\ref{Extension for any strings from status equivalence to syntactic congruence}, 
we have \(x\sim_L\rho(x)\) and \(y\sim_L\rho(y)\). 
By transitivity of \(\sim_L\), we obtain \(x\sim_L y\). This proves \((**)\).  \hfill $\square$

\begin{theorem}[Correctness of the \texttt{Promote\_to\_P} Strategy]
\label{theorem-correctness-of-the-promote_to_p-strategy}
Let \(OT=(I,P,S,T)\) be a closed, separable, and right-consistent
observation table maintained by the \texttt{Promote\_to\_P} strategy,
so that \(I\subseteq P\).
Let \(\mathcal M\) be its hypothesis acceptor, and assume that
\(\mathcal M\) passes the equivalence query, that is,
\(L(\mathcal M)=U\).
Then, for all \(x,y\in\Sigma^*\),
\[
x\equiv_{\mathrm{status}}y
\Longleftrightarrow x\sim_U y.
\]
\end{theorem}
\noindent\textbf{Proof.}
The statement is similar to Theorem~\ref{theorem-correctness-of-extended-lstar-algorithm-for-post-check-strategy} and we
only need to show that \(\equiv_{\mathrm{status}}\) implies \(\sim_U\).

First, we show that the table is monoid consistent because in this strategy, there is no monoid consistency check. 
Let \(t\in I\Sigma\) and \(v\in I\) satisfy
\(t\equiv_{\mathrm{status}}v\).
Since \(I\subseteq P\), equality of their full rows gives
\(T(i,t,s)=T(i,v,s)\) for every \(i\in I\) and \(s\in S\).
By the definition of \(T\), this means
\(\chi(its)=\chi(ivs)\).
Thus monoid consistency holds.

The table therefore satisfies monoid consistency, so both
Lemma~\ref{Extension for each transition step from status equivalence to syntactic congruence}
and
Lemma~\ref{Extension for any strings from status equivalence to syntactic congruence}
apply. The remainder of the proof follows the same argument as in
Theorem~\ref{theorem-correctness-of-extended-lstar-algorithm-for-post-check-strategy}.
Hence status equivalence coincides with the syntactic congruence
of \(U\), and the returned hypothesis is isomorphic to the right
Cayley graph of its syntactic monoid, equipped with its initial
and accepting states.
\hfill $\square$

\section{Complexity Analysis}\label{complexity-analysis}

We distinguish membership-query complexity, equivalence-query
complexity, and internal computation. The analysis below treats the
two strategies separately and counts repeated checks separately from
distinct queries to the teacher.

\subsection{Cost Model and Common Operations}\label{cost-model-common-operations}

Let \(k=|\Sigma|\ge1\), and let \(n,p,s\) denote the final sizes of
\(I,P,S\) in a completed run. Since these sets only grow, their final
sizes also bound their sizes at every earlier stage. Let \(M\) be the
size of the target syntactic monoid and \(m\ge1\) an upper bound on
counterexample length. Separability gives \(n\le M\): syntactically
congruent words cannot have different observation rows.

We analyze the variant that adds one distinguishing suffix for each
counterexample or right-consistency refinement. Membership results are
cached throughout the run. Internal bounds count table access, row
comparison, hashing, and the operations described below; they exclude
oracle execution time, string copying, and ordering the sets for
presentation. Hash-table lookups have expected constant cost after
computing a key's hash. Thus these are bounds in a row-operation model,
rather than full bit-complexity bounds.

\textbf{Observation Table and Membership Queries.}
A full row has \(ps\) Boolean entries. Since
\(|I\cup I\Sigma|\le(k+1)n\), the table has \(O(knps)\) cells.
Filling all cells from scratch requires at most that many MQs.
With a persistent cache, repeated table filling does not multiply the
number of distinct MQs: every queried table word belongs to the final
set \(P(I\cup I\Sigma)S\). Different cells may also represent the same
word, so the cell count is an upper bound, not an exact query count.

\textbf{Closedness and Separability.}
Hash the rows in \(I\), then check the at most \(kn\) boundary rows
against them. Computing a full-row hash takes \(O(ps)\), giving
\(O(knps)\) expected work for a complete pass, including row
construction if necessary. Adding only previously unseen rows
preserves separability without a separate all-pairs check.

\textbf{Right Consistency.}
The direct implementation considers \(O(n^2)\) pairs of
representatives. For each suffix-equivalent pair it checks \(k\)
letters, comparing suffix rows of length \(s\). Hence a complete
pass costs \(O(kn^2s)\), including the worst-case scan for a
distinguishing suffix. These membership values already occur in the
table at left context \(\epsilon\in P\); the check therefore adds
no distinct MQs beyond table filling. A refinement adds \(as\) to
\(S\), after which newly required table cells must be filled.

\textbf{Hypothesis Construction.}
Hash-based state and transition construction costs \(O(knps)\).
If distinct full rows are also sorted, as in the implementation,
sorting contributes \(O(nps\log(2+n))\). We write
\(H=knps+nps\log(2+n)\) for this per-construction upper bound.

\textbf{Counterexample Processing.}
Binary-search refinement~\cite{rivest1993inference} uses
\(O(\log(2+m))\) MQs per counterexample. Tracing the hypothesis along
the counterexample takes \(O(m)\) internal transition steps.
The subsequent table expansion is counted under table filling.
If \(f_E\) and \(f_R\) count counterexamples and right-consistency
refinements, respectively, the single-suffix variant satisfies
\(s\le1+f_E+f_R\). In particular, growth of \(S\) cannot be charged
only to failed equivalence queries.

For cumulative internal bounds, let \(r_T\) count table-filling and
closedness passes, \(r_R\) right-consistency passes, and \(r_H\)
hypothesis constructions, including checks that succeed. Define
\[
B=r_Tknps+r_Rkn^2s+r_HH+f_Em.
\]
Each strategy uses its own final sizes and pass counts in this
expression. Let \(q_E\) denote its actual number of teacher EQ calls.
We retain these counts explicitly instead of assuming an unproved
bound on the number of learning rounds.

\subsection{The \texttt{post\_check} Strategy}\label{complexity-post-check}

\textbf{Table Size.}
In the variant of Section~\ref{strategy-post-check}, each new left
context is a representative from \(I\). Thus \(P\subseteq I\) and
\(p\le n\le M\). The table bound remains \(O(knps)\), retaining
the benefit of using only the left contexts actually needed.

\textbf{Monoid Consistency Checks.}
Closedness and separability determine one representative \(v\in I\)
for each \(t\in I\Sigma\), so all \(n\) possible representatives
need not be tried. There are at most \(kn\) boundary pairs, and each
is tested with at most \(ns\) contexts \((i,s')\in I\times S\).
Thus the probes cost \(O(kn^2s)\) per complete check. Forming or
hashing the row-to-representative map adds at most \(O(knps)\).
Tests with \(i\in P\) can be skipped, since their results already
agree by full-row equality.

Let \(q_M\) be the number of monoid checks, and \(f_P\) their number
of failures. Each failure adds a previously absent word to \(P\),
and the final successful check terminates the run. Consequently,
\(q_M=f_P+1\le p\). The cumulative internal cost is bounded by
\[
C_{\mathrm{post}}=O\bigl(B+q_M(knps+kn^2s)\bigr).
\]

\textbf{Distinct Membership Queries.}
All monoid probes over the entire run have the form \(its'\) or
\(ivs'\), with \(i,v\in I\), \(t\in I\Sigma\), and \(s'\in S\).
Their union therefore contains at most \(O(kn^2s)\) words.
Persistent caching makes this a whole-run MQ bound, without a
factor of \(q_M\). Together with table filling and counterexamples,
\[
Q_{\mathrm{post}}=O\bigl(knps+kn^2s+f_E\log(2+m)\bigr).
\]
Repeated probes still contribute to internal computation even when
their MQ results are cached.

\textbf{Equivalence Queries.}
The EQ cost is \(q_E\) teacher calls. A successful language EQ can
be followed by further prefix refinements, so language confirmation
and completion of monoid learning are separate events. If a
confirmed language equivalence is reused, skipped checks do not
count as new teacher calls. No bound \(q_E\le M\) is assumed here.

\subsection{The \texttt{Promote\_to\_P} Strategy}\label{complexity-promote-to-p}

\textbf{Table Size and Promotion.}
Starting with \(I=P=\{\epsilon\}\) and adding every new representative
to both sets gives \(P=I\), provided counterexample processing only
refines \(S\). Hence \(p=n\), and the table has \(O(kn^2s)\) cells.
If another counterexample handler inserts additional left contexts,
only \(I\subseteq P\) is guaranteed; in that case the more general
bound \(O(knps)\) must be retained.

\textbf{Internal Checks.}
For the single-suffix variant, closedness costs \(O(kn^2s)\) per
pass, and right consistency also costs \(O(kn^2s)\) per pass.
The hypothesis-construction bound is \(H\) with \(p=n\).
Since all representatives are already left contexts, monoid
consistency is included in full-row equality and \(q_M=0\).
Thus the cumulative internal cost is
\[
C_{\mathrm{promote}}=O(B),\qquad p=n.
\]

\textbf{Distinct Membership Queries and EQs.}
No additional monoid-probe queries are needed. Therefore,
\[
Q_{\mathrm{promote}}=O\bigl(kn^2s+f_E\log(2+m)\bigr).
\]
The strategy stops at its first successful EQ, so
\(q_E=f_E+1\). Right-consistency refinements take place before
submitting a hypothesis and may change how many counterexamples
are needed; they must still be included in the internal cost.

\subsection{Comparison of the Two Strategies}\label{complexity-strategy-comparison}

The main trade-off is between table width and explicit monoid
checks. For the same \(n,s\), \texttt{post\_check} uses \(ps\) tests
per row, while \texttt{Promote\_to\_P} uses \(ns\) in the
single-suffix variant. Conversely, only \texttt{post\_check} pays
for the repeated explicit monoid scans. Both strategies pay for
right consistency; avoiding monoid checks does not remove this cost.

Their cached MQ bounds can have the same worst-case order even
when internal costs differ substantially. The values of \(p,s\),
the number of repeated passes, and oracle costs determine which
strategy is cheaper on a particular target. Fewer EQs for one
strategy are not guaranteed solely by these upper bounds.

The inequality \(p+s\le M\) is not used: in the promotion strategy,
\(p=M\) at successful completion while \(s\ge1\). Likewise, a
bound on \(s\) or on the number of passes requires a separate
progress argument that includes right-consistency refinements.
The present bounds therefore keep these parameters explicit rather
than assert an unsupported total complexity solely in \(M\).

\section{Conclusion}\label{conclusion}

This report presented an observation-table-based approach for learning
the syntactic monoid of an unknown regular language. The main components
of the method are the construction of a hypothesis acceptor, the use of
membership and equivalence queries, and the additional monoid
consistency check. When the observation table is closed, separable,
right consistent, and monoid consistent, the resulting acceptor
represents the syntactic monoid recognizing the unknown regular language
\(U\).


\bibliographystyle{ieicetr}
\bibliography{references}
\end{document}
